% =================================================================
% CAPÍTULO 4: AVALIAÇÃO DE RESULTADOS (TEMPLATE)
% Objetivo: apresentar resultados e como foram avaliados.
% Dicas:
%  - Prefira tabelas LaTeX (booktabs) a imagens de tabelas.
%  - Referencie figuras/tabelas no texto (ex.: ver Tabela~\ref{tab:resultado-exemplo}).
% =================================================================
\section{Avaliação de Resultados}

[Apresente os resultados e como foram avaliados, relacionando-os aos objetivos.]

\subsection{Métricas e Procedimentos}
[Liste métricas utilizadas, procedimentos de avaliação e critérios de aceitação.]

\subsection{Tabela de Exemplo}
\begin{table}[h]
  \centering
  \caption{Exemplo de resultados comparativos}
  \label{tab:resultado-exemplo}
  \begin{tabular}{lcc}
    \toprule
    Métrica & Antes & Depois \\
    \midrule
    Exemplo 1 & X & Y \\
    Exemplo 2 & X & Y \\
    \bottomrule
  \end{tabular}
\end{table}

\subsection{Figura de Exemplo}
\begin{figure}[h]
  \centering
  \includegraphics[width=0.85\linewidth]{template_fig1.jpg}
  \caption{[Legenda da figura de resultados]}
  \label{fig:resultados-exemplo}
\end{figure}
